\chapter{アルゴリズムと数式}
\label{chap:algorithm}
\chapterauthor{\AuthorTwo}

\section{総和を数式で表す}
第\ref{chap:python}章の処理は、インライン数式 $S_n=1+2+\cdots+n$ で表せます。
\begin{definition}[三角数]
  正の整数 $n$ に対し、$T_n=\sum_{k=1}^{n}k$ を第$n$三角数と呼びます。
\end{definition}
\begin{theorem}
  三角数は次の式で求められます。
  \begin{equation}
    T_n=\frac{n(n+1)}{2}.
    \label{eq:triangular}
  \end{equation}
\end{theorem}
\begin{proof}
  同じ和を逆順にも並べて加えると、各組の和は $n+1$ です。
  \begin{align}
    2T_n &= (1+n)+(2+(n-1))+\cdots+(n+1) \\
         &= n(n+1).
  \end{align}
  両辺を2で割れば式\eqref{eq:triangular}を得ます。
\end{proof}
番号を付けない別行立て数式も利用できます。
\[
  T_5=\frac{5\cdot6}{2}=15
\]

\section{処理の流れを図にする}
図\ref{fig:flow}はLaTeXの枠と矢印だけで作った図です。
\begin{figure}[htbp]
  \centering
  \fbox{入力：$n$}\enspace$\longrightarrow$\enspace
  \fbox{計算：$n(n+1)/2$}\enspace$\longrightarrow$\enspace
  \fbox{出力：$T_n$}
  \caption{三角数を求める処理の流れ}
  \label{fig:flow}
\end{figure}

\subsection{参考資料を示す}
数式を説明するときも、出典を明示しましょう。
本書の日本語組版にはjlreq\cite{ref:jlreq}を使用しています。
\begin{quote}
  ここには引用文を記載し、出典と該当ページを添えます。
  この文章自体は引用の書式を示すための例文です。
\end{quote}
